\documentclass[english,twoside, openright, hidelinks, 11 pt, class=report,crop=false]{standalone}
\input{../bm}
\input{../../preamb}


\begin{document}
\subsection*{Kapittel \ref{Funk}}

\opr{opginvfindtoppbunn}
\abch{
	\item Bunnpunkt $ \left(e^{-\frac{1}{2}}, -2e^{-1}\right)$
	\item Bunnpunkt $ \left(\frac{1}{2}, 0\right) $ og toppunkt $\left(-\frac{3}{2}, 8e^{-\frac{3}{2}}\right)$
}

\opr{opgfunkekstrm}
$\frac{b+c}{2}$

\opr{1tv21d1opg7}
$a=20$

\opr{opgfunkfinnf}
\abch{
	\item $f(x)=-\frac{5}{28}(x-3)(x+4)$
	\item $f(x)=-\frac{-25}{9}(x-10)(x+1)$
	\item $f(x)=-(x-8)(x-10)+9$
} 

\opr{opginvkonvekskonkav}
\selos

\opr{opgfunkasymp}
\abch{
	\item Vertikal: $x=2$, horisontal: $y=0$, skjæringspunkt: $(0, -2)$
	\item Vertikal: $x=-3$, horisontal: $y=0$, skjæringspunkt: $(0, \frac{3}{7})$
	\item Vertikal: $x=\pm4$, horisontal: $y=1$, skjæringspunkt: $(0, 0)$
	\item Vertikal: $x=2$, horisontal: $y=1$, skjæringspunkt: $(0, -\frac{3}{2})$
	\item Vertikal: $x=0$, horisontal: $y=1$, skjæringspunkt: $(0, -8)$
	\item Vertikal: $x=-\frac{1}{2}$, horisontal: $y=6$, skjæringspunkt: $(0, -3)$
}

\opr{1tv22d1opg5}
$f(x)=\frac{3}{x+2}+3 $ eller $f(x)=\frac{3x+9}{x+2}$ (Obs! Flere mulige svar)

\opr{opgfunkomv}
\abch{
	\item $ g(y)=\frac{y}{3} $ 
	\item $ g(y)=\frac{2-y}{9} $
	\item $ g(y)=\frac{2y+14}{5} $
	\item $ g(y)=\frac{y}{3}+5 $ 
	\item $ g(y)=y^2 $
	\item $ g(y)=y^3 $
	\item $ g(y)=(y-9)^{\frac{1}{4}} $
}

\opr{opgfunkinvlogarithm}
\abch{
	\item $g(y)=\ln(y-2)$
	\item $g(y)=e^y-5$
	\item $g(y) = e^{\frac{1}{y}}$
}

\opr{r1v26d1opg6}
$g'(1)=e$

\opr{r1h22d1opg1}	
Avgjør om funksjonen har en omvendt funksjon for hele definisjonsmengden. Begrunn svaret ditt.
\abc{
	\item Nei. Ikke strengt voksende på hele intervallet.
	\item Ja. Strengt avtagende på hele intervallet.
}

\opr{opgfunkinvfinna}
$a=-1$


\opr{opgfunkdrtppktbnpkt}
Den er strengt voksende eller strengt avtagende på intervallet.

\grubr{opginvwhich}
Endepunkt og kritiske punkt.

\grubr{1th24d2opg5} 
$a=-\frac{1}{4}$, $b=\frac{3}{2}$, $c=-3$, $d=10$

\grubr{1tv23d1opg4}
$f(x)=\frac{3x-6}{x-1}$

\grubr{1th24d2opg3}
$f(x)=\frac{4x+12}{x-2}$ (Obs! Flere mulige svar)

\grubr{1tv26d1opg8}
\abch{
	\item \selos
	\item \selos
}

\grubr{1tv25opg4}
\abch{
	\item $x\in\{-2, 1, 8\}$
	\item c
}

\grubr{1th23d1opg5}
\abch{
	\item c
	\item f
}

\grubr{opgfunkdrviskonv}
\selos

\grubr{r1v22d2opg3}
$I=[-\sqrt{2}, 1]$


\grubr{r1v26d1opg5}{R1V26}
\abch{
	\item Ikke riktig (\selos).
	\item Riktig (\selos).
}

\grubr{r1h23d2opg2}
\abc{
	\item \selos
	\item $k=1$
	\item $k\leq1$
}

\grubr{opgderrectincircle}
\selos


\grubr{r1v24d2opg6}
$C = (1, 3)$

\grubr{r1v22d2opg5}
\abch{
	\item $\frac{9}{32}$.
	\item $a=\frac{1}{2}$.
}

\grubr{r1h24d2opg5ish}
\abch{
	\item $f$
	\item $f(x)=x^2+3$, $g(x)= \left\lbrace{
		\begin{array}{rcr}
			x-1 &,&x<1 \br
			2x-4   &,& x\geq 1
		\end{array}
	}\right.$, $h=(x-1)(x+1)(x+2)$
	\item \selos
}

\grubr{r1v23d2opg3}
1. Usann, 2. Sann, 3. Sann.


\grubr{r1v22d2opg8}
\selos

\grubr{R1V23D2opg4}
\abch{
	\item $f$: nei, $g$: ja, $h$: ja, $k$: ja.
	\item $D_g=[-4, 4]$, $D_h=[-5, -2)\cup[-2, 4]$, $D_k=[-4, 4]$.
}

\grubr{r1v22d2opg6}
$s=3$, 16

\grubr{r1v23d2opg6}
\abch{
	\item $5$
	\item $\frac{3\sqrt{5}}{5}$
}


\grubr{r1h25d2opg6}
\abch{
	\item Funksjon og andrederivert: A og G, E og H, B og C, F og D, 
	\item A, B, C og G fordi de er injektive
}

\grubr{r1h23d2opg4} 
\abch{
	\item 100
	\item 64
	\item 120
}


\grubr{opginvsym2}
\abch{
	\item $B=(b, a)$
	\item $\vec{r}=[b-a, b-a]$, $45^\circ$
}

\grubr{opginvpointb}
\abch{
	\item $\vec{v}=t[1, 2-\sqrt{3}]$.
} \\
\abchs{2}{
	\item $f(t)=(8-4\sqrt{3})t^2+(7-5\sqrt{3})t - 12(13-7\sqrt{3})$ (Obs! Flere uttrykk mulig)
	\item \selos
	\item $B=(4, 4(2-\sqrt{3}))$
}

\grubr{opginvareaoftriangle}
\abch{
	\item $A=\left(\frac{\sqrt{3}}{4}, 0\right)$, $B=\left(\frac{\sqrt{3}}{2}, \frac{4}{3}\right)$
	\item \selos
	\item $\frac{\sqrt{3}}{8}\left(1+3\sqrt{3}\right)$
	\item $\frac{1}{2}\sqrt{10\sqrt{3}+41}$
	\item $(x-2)(x^2-2x-5)$
	\item $(1, -6)$
	\item $x=\sqrt{6}+1\vee x=1-\sqrt{6}$
}

\grubr{opginvsym}
\selos

\grubr{opginvaxba}
\abch{
\item \selos
\item $g(a)=\sqrt{a^2+2}$
}

\grubr{opginvconvx}
\selos


\end{document}